%% SYNTHETIC TEST UNIT generated by writing/tools/make_outline_book.py from FINAL_OUTLINE_v2 -- not book text
\chapter{其他共线过程：单举湮灭、SIDIS、光生、极化与横度}\label{ch:coll}
本单元为测试用的合成文本（由 FINAL\_OUTLINE\_v2 自动生成），陈述内容从略。

\section{第 1 节（测试）}\label{sec:coll-sia}

\begin{definition}[单举湮灭的运动学]\label{def:coll-sia}
（测试用合成陈述：内容从略。）
\uses{def:pdf-ff}
\end{definition}

\begin{lemma}[全部为末态的共线子图的极点在同一半平面 \tagR]\label{lem:coll-final-poles}
（$R_{\mathrm{PT}}$：在重整化微扰论中对红外正规化的部分子态严格。）
（测试用合成陈述：内容从略。）
\uses{lem:appO-lcopt, def:coll-sia}
\end{lemma}

\begin{lemma}[未观测喷注与软辐射的幺正性相消 \tagA]\label{lem:coll-unobserved-jet}
（测试用合成陈述：内容从略。）
\uses{prop:method-fsi, prop:prelim-unitarity}
\end{lemma}
\begin{proofsketch}测试。\end{proofsketch}
\paragraph{地位与缺口}测试：本陈述的地位由\cref{asm:appN-cutsum}决定。

\begin{proposition}[SIA 的领头区域 \tagA]\label{prop:coll-sia-regions}
（测试用合成陈述：内容从略。）
\uses{thm:ir-leading-regions, lem:coll-final-poles, lem:coll-unobserved-jet}
\end{proposition}
\begin{proofsketch}测试。\end{proofsketch}
\paragraph{地位与缺口}测试：本陈述的地位由\cref{asm:appN-envelope}、\cref{asm:appN-cutsum}决定。

\begin{theorem}[$e^+e^-\to hX$ 的因子化 \tagA]\label{thm:coll-sia}
（测试用合成陈述：内容从略。）
非微扰地位：\cref{asm:intro-bridge}。
\uses{thm:method-template, prop:coll-sia-regions, lem:coll-final-poles, lem:coll-unobserved-jet, def:pdf-ff, rem:pdf-ff-multren, asm:appN-multren, asm:intro-bridge}
\end{theorem}
\begin{proofsketch}测试。\end{proofsketch}
\paragraph{地位与缺口}测试：本陈述的地位由\cref{asm:appN}、\cref{rem:pdf-ff-multren}、\cref{asm:appN-multren}决定。

\section{第 2 节（测试）}\label{sec:coll-sidis}

\begin{remark}[系数函数与检验]\label{rem:coll-sia-nlo}
（测试用合成陈述：内容从略。）
% uses: none
\end{remark}

\begin{example}[领头阶 $e^+e^-\to hX$]\label{ex:coll-sia-lo}
（测试用合成陈述：内容从略。）
\uses{thm:coll-sia, def:coll-sia}
\end{example}

\begin{theorem}[共线 SIDIS 的因子化 \tagA]\label{thm:coll-sidis}
（测试用合成陈述：内容从略。）
非微扰地位：\cref{asm:intro-bridge}。
\uses{thm:method-template, thm:disall-factorization, thm:coll-sia, lem:coll-final-poles, asm:intro-bridge}
\end{theorem}
\begin{proofsketch}测试。\end{proofsketch}
\paragraph{地位与缺口}测试：本陈述的地位由\cref{thm:disall-factorization}、\cref{thm:coll-sia}决定。

\begin{theorem}[直接光生与光子结构函数 \tagA]\label{thm:coll-direct-photo}
（测试用合成陈述：内容从略。）
非微扰地位：\cref{asm:intro-bridge}。
\uses{thm:disall-factorization, asm:intro-bridge}
\end{theorem}
\begin{proofsketch}测试。\end{proofsketch}
\paragraph{地位与缺口}测试：本陈述的地位由\cref{thm:disall-factorization}决定。

\begin{proposition}[光子结构函数的点状部分 \tagR]\label{prop:coll-photon-sf}
（$R_{\mathrm{OPE}}$：以 OPE 为前提而严格。QCD 中含夸克的微扰 OPE 是全阶论证（\cref{thm:dis-ope}），严格证明只对标量理论与纯 Yang--Mills 的规范不变算符给出；对强子态还须非微扰 OPE（\cref{asm:intro-np-ope}）。）
（测试用合成陈述：内容从略。）
非微扰地位：\cref{asm:intro-wightman}、\cref{asm:intro-composite}、\cref{asm:intro-np-ope}、\cref{asm:intro-af-coeff}。
\uses{thm:dis-ope, prop:dis-ope-rg, asm:intro-wightman, asm:intro-composite, asm:intro-np-ope, asm:intro-af-coeff}
\end{proposition}

\section{第 3 节（测试）}\label{sec:coll-photo}

\begin{conjecture}[分辨光生 \tagC]\label{conj:coll-resolved-photo}
（测试用合成陈述：内容从略。）
\uses{thm:coll-direct-photo}
\end{conjecture}
\paragraph{地位与缺口}测试：本陈述的地位由\cref{conj:dy-coloured-final}决定。

\begin{proposition}[衍射光生中的因子化破坏 \tagV]\label{prop:coll-diffractive-photo}
（测试用合成陈述：内容从略。）
\uses{conj:coll-resolved-photo}
\end{proposition}
\paragraph{地位与缺口}测试：本陈述的地位由\cref{thm:break-ddis}决定。

\begin{proposition}[质子的光子分布 \tagA]\label{prop:coll-luxqed}
（测试用合成陈述：内容从略。）
非微扰地位：\cref{asm:intro-wightman}、\cref{asm:intro-composite}。
\uses{def:pdf-gluon, def:dis-kinematics, asm:intro-wightman, asm:intro-composite}
\end{proposition}
\begin{proofsketch}测试。\end{proofsketch}
\paragraph{地位与缺口}测试：见文献。

\begin{definition}[极化结构函数]\label{def:coll-g1}
（测试用合成陈述：内容从略。）
\uses{def:dis-kinematics, def:prelim-gamma5}
\end{definition}

\begin{theorem}[$g_1$ 的奇数矩 \tagR]\label{thm:coll-g1-moments}
（$R_{\mathrm{OPE}}$：以 OPE 为前提。）
（测试用合成陈述：内容从略。）
非微扰地位：\cref{asm:intro-wightman}、\cref{asm:intro-composite}、\cref{asm:intro-np-ope}、\cref{asm:intro-af-coeff}、\cref{asm:intro-growth}。
\uses{thm:dis-ope, def:coll-g1, thm:dis-bjorken, prop:appI-larin, asm:intro-wightman, asm:intro-composite, asm:intro-np-ope, asm:intro-af-coeff}
\end{theorem}

\section{第 4 节（测试）}\label{sec:coll-polarized}

\begin{proposition}[一阶矩的方案依赖 \tagR]\label{prop:coll-anomaly-scheme}
（$R_{\mathrm{PT}}$：在重整化微扰论中对红外正规化的部分子态严格。）
（测试用合成陈述：内容从略。）
\uses{def:appI-hvbm, prop:appI-larin, def:coll-g1}
\end{proposition}

\begin{proposition}[Ellis–Jaffe 求和规则作为预言 \tagV]\label{prop:coll-ellis-jaffe}
（测试用合成陈述：内容从略。）
\uses{thm:coll-g1-moments}
\end{proposition}
\paragraph{地位与缺口}测试：见文献。

\begin{conjecture}[Burkhardt–Cottingham 求和规则 \tagC]\label{conj:coll-bc}
（测试用合成陈述：内容从略。）
\uses{def:coll-g1}
\end{conjecture}
\paragraph{地位与缺口}测试：本陈述的地位由\cref{asm:intro-growth}、\cref{prop:dis-low-moments}决定。

\begin{proposition}[Wandzura–Wilczek 关系的一般失效 \tagV]\label{prop:coll-ww}
（测试用合成陈述：内容从略。）
\uses{def:coll-g1}
\end{proposition}
\paragraph{地位与缺口}测试：见文献。

\begin{definition}[横度分布]\label{def:coll-transversity}
（测试用合成陈述：内容从略。）
\uses{def:pdf-polarized}
\end{definition}

\section{第 5 节（测试）}\label{sec:coll-transversity}

\begin{proposition}[横度的演化 \tagR]\label{prop:coll-transversity-evolution}
（$R_{\mathrm{def}}$：由算符定义直接推出。）
（测试用合成陈述：内容从略。）
非微扰地位：\cref{asm:intro-wightman}、\cref{asm:intro-composite}。
\uses{def:coll-transversity, def:pdf-polarized, asm:intro-wightman, asm:intro-composite}
\end{proposition}

\begin{proposition}[Soffer 不等式（LO 演化的保持性） \tagR]\label{prop:coll-soffer}
（$R_{\mathrm{PT}}$：在重整化微扰论中对红外正规化的部分子态严格。）
（测试用合成陈述：内容从略。）
\uses{prop:coll-transversity-evolution, prop:pdf-lo-kernels}
\end{proposition}

\begin{remark}[单举 DIS 中无单靶自旋横向不对称]\label{rem:coll-christ-lee}
（测试用合成陈述：内容从略。）
% uses: none
\end{remark}
